% Generated draft: requires mathematical and source review.
\documentclass[11pt]{article}
\usepackage[margin=25mm]{geometry}
\usepackage{fontspec}
\setmainfont{Latin Modern Roman}
\usepackage{amsmath,amssymb,mathtools,unicode-math}
\setmathfont{Latin Modern Math}
\usepackage{enumitem,needspace,xurl,hyperref}
\setlength{\parindent}{0pt}
\setlength{\parskip}{5pt}
\setlength{\emergencystretch}{3em}
\newcommand{\N}{\mathbb{N}}
\newcommand{\Z}{\mathbb{Z}}
\newcommand{\Q}{\mathbb{Q}}
\newcommand{\R}{\mathbb{R}}
\newcommand{\C}{\mathbb{C}}
\newcommand{\F}{\mathbb{F}}
\newcommand{\A}{\mathbb{A}}
\newcommand{\PP}{\mathbb{P}}
\newcommand{\E}{\mathbb{E}}
\newcommand{\eps}{\varepsilon}
\newcommand{\dd}{\,\mathrm{d}}
\begin{document}
\section*{UnsolvedMath ID: 800013}
\textbf{Source problem code:} catalogue entry
\subsection{Definitions and mathematical statement}
A semigroup is a nonempty set equipped with an associative binary operation. For a semigroup \(S\), let
\[
\operatorname{End}(S)=\{\varphi:S\to S: \varphi(xy)=\varphi(x)\varphi(y)\text{ for all }x,y\in S\}.
\]
A variety of semigroups is a class defined by identities in the signature consisting of one binary operation. For a class \(K\) of semigroups and a positive integer \(n\), put
\[
\mathcal F_K(n)=\{S\in K:S\text{ is finite and }|S|\ge n\}.
\]
Whenever \(\mathcal F_K(n)\ne\varnothing\), define its endomorphic function at \(n\) by
\[
f_K(n)=\min\{|\operatorname{End}(S)|:S\in\mathcal F_K(n)\}.
\]
The source calls \(K\) \emph{endomorphic} when \(f_K\) is unbounded.

\textbf{Problem (S. Kublanovsky).} Describe the endomorphic varieties of semigroups. For any variety \(V\) having finite members of arbitrarily large cardinality, this asks for a characterization of the condition
\[
\forall B\in\Z_{\ge0}\ \exists n\in\Z_{\ge1}\ \forall S\in V,\qquad
\bigl(S\text{ finite and }|S|\ge n\bigr)
\Longrightarrow |\operatorname{End}(S)|>B.
\]
Under this hypothesis all the defining minima exist, and the displayed condition is exactly unboundedness of \(f_V\).

\textbf{Unresolved domain qualification.} The source defines \(f_K(n)\) for every natural number without specifying what happens when \(\mathcal F_K(n)\) is empty. Thus the supplied formulation does not determine whether varieties without arbitrarily large finite members are excluded, or how they should be treated. The quantified formulation above describes the target only when the stated existence hypothesis holds; it does not settle this omitted convention. The source also reports that some sufficient conditions were obtained in the reference labelled \texttt{[kubl]}, whose bibliographic details are not supplied.
\subsection{Short English statement}
Classify varieties of semigroups for which the minimum number of endomorphisms of a finite member with at least \(n\) elements is unbounded as \(n\) grows. This definition requires clarification when the variety has no sufficiently large finite members. The problem is attributed to S. Kublanovsky; the supplied registry labels it open.
\subsection{Sources}
\begin{itemize}
\item Supplied local UnsolvedMath registry record, permanent ID \(800013\), exact source problem code \texttt{catalogue entry}, in the AMR Open Problem Lists collection. The registry lists Bershadsky, Kublanovsky, and Mashevitzky under \texttt{proposed\_by}, with year \(2026\), and labels the record \texttt{open}; these are supplied metadata, not independently verified conclusions.
\item Bershadsky, Kublanovsky, and Mashevitzky, \emph{The Lyapin's Notebook: Unsolved Problems in Semigroup Theory} (2026), Problem 13. Supplied source URL: \url{https://arxiv.org/abs/2604.04763}. The background records access on 2026-10-16 and extraction by \texttt{arxiv\_latex}; these are supplied provenance details, not a claim of independent source access or verification. The problem text attributes the question specifically to S. Kublanovsky.
\item The supplied statement cites \texttt{[kubl]} for some sufficient conditions. No bibliographic entry, title, date, or link for this reference is included in the supplied record.
\end{itemize}
\end{document}
